% Bagian: turing-machines
% Bab: undecidability
% Subbagian: unsolvability-decision-problem

\documentclass[../../../include/open-logic-section]{subfiles}

\begin{document}

\olfileid[id]{tur}{und}{uns}
\olsection{Masalah Keputusan Tidak Dapat Diselesaikan}

\begin{thm}
\ollabel{thm:decision-prob}
Masalah keputusan tidak dapat diselesaikan: tidak ada mesin Turing~$D$ yang,
ketika dijalankan dengan pita yang memuat !!a{sentence} logika orde pertama
$!B$ sebagai masukan, pada akhirnya berhenti dan menghasilkan~$1$ jika dan
hanya jika $!B$ valid, serta menghasilkan~$0$ jika tidak.
\end{thm}

\begin{proof}
Andaikan masalah keputusan dapat diselesaikan, yakni andaikan terdapat mesin
Turing~$D$. Maka kita dapat menyelesaikan masalah penghentian sebagai berikut.
Kita mengonstruksi mesin Turing~$E$ yang, ketika diberi nomor~$e$ dari mesin
Turing~$M_e$ dan masukan~$w$, menghitung !!{sentence} yang bersesuaian
$!T(M_e, w) \lif !E(M_e, w)$, lalu berhenti sambil memindai petak paling kiri
pada pita. Mesin $E \concat D$, ketika diberi masukan $e$ dan~$w$, mula-mula
akan menghitung $!T(M_e, w) \lif !E(M_e, w)$ lalu menjalankan mesin masalah
keputusan~$D$ pada masukan tersebut. $D$ berhenti dengan keluaran~$1$ jika dan
hanya jika $!T(M_e, w) \lif !E(M_e, w)$ valid, serta menghasilkan~$0$ jika
tidak. Menurut \olref[ver]{lem:halt-if-valid} dan
\olref[ver]{lem:valid-if-halt}, $!T(M_e, w) \lif !E(M_e, w)$ valid jika dan
hanya jika $M_e$ berhenti pada masukan~$w$. Jadi, ketika diberi masukan $e$
dan~$w$, $E\concat D$ berhenti dengan keluaran~$1$ jika dan hanya jika $M_e$
berhenti pada masukan~$w$, serta berhenti dengan keluaran~$0$ jika tidak.
Dengan kata lain, $E \concat D$ akan menyelesaikan masalah penghentian.
Namun, menurut \olref[hal]{thm:halting-problem}, kita mengetahui bahwa mesin
Turing semacam itu tidak mungkin ada.
\end{proof}

\begin{cor}\ollabel{cor:undecidable-sat}%
Tidak dapat diputuskan apakah suatu !!{sentence} sebarang logika orde pertama
dapat dipenuhi.
\end{cor}

\begin{proof}
  Andaikan keterpenuhan dapat diputuskan oleh mesin Turing~$S$. Maka masalah
  keputusan dapat kita selesaikan sebagai berikut. Jika diberi
  !!a{sentence}~$B$ sebagai masukan, pindahkan $!B$ satu petak ke kanan.
  Kembalilah ke petak~$1$ lalu tuliskan simbol~$\lnot$.

  Sekarang jalankan mesin Turing~$S$. Mesin itu pada akhirnya berhenti dengan
  keluaran $1$ (jika $\lnot !B$ dapat dipenuhi) atau~$0$ (jika $\lnot !B$
  tidak dapat dipenuhi) pada pita. Jika terdapat~$\TMstroke$ pada petak~$1$,
  hapuslah; jika petak~$1$ kosong, tuliskan~$\TMstroke$, kemudian berhenti.

  Mesin Turing ini selalu berhenti, dan keluarannya adalah~$1$ jika dan hanya
  jika $\lnot !B$ tidak dapat dipenuhi, serta~$0$ jika tidak. Karena $!B$
  valid jika dan hanya jika $\lnot !B$ tidak dapat dipenuhi, mesin itu
  menghasilkan~$1$ jika dan hanya jika $!B$ valid, serta~$0$ jika tidak;
  yakni, mesin tersebut akan menyelesaikan masalah keputusan.
\end{proof}

\begin{explain}
Jadi, tidak ada mesin Turing yang selalu memberikan jawaban ``ya'' atau
``tidak'' yang benar atas pertanyaan ``Apakah $!B$ merupakan !!{sentence}
logika orde pertama yang valid?'' Akan tetapi, memang \emph{ada} mesin Turing
yang selalu memberikan jawaban ``ya'' yang benar---tetapi tidak berhenti jika
jawabannya ``tidak.'' Hal ini mengikuti teorema kesahihan dan kelengkapan
logika orde pertama serta fakta bahwa !!{derivation} dapat dienumerasi secara
efektif.
\end{explain}

\begin{thm}
  \ollabel{thm:valid-ce}%
  Validitas !!{sentence} orde pertama bersifat semidapat diputuskan: terdapat
  mesin Turing~$E$ yang, ketika dijalankan dengan pita yang memuat
  !!a{sentence} logika orde pertama~$!B$ sebagai masukan, pada akhirnya
  berhenti dan menghasilkan~$1$ jika dan hanya jika $!B$ valid, tetapi tidak
  berhenti jika tidak.
\end{thm}

\begin{proof}
  Semua !!{derivation} yang mungkin dalam logika orde pertama dapat dihasilkan
  satu demi satu oleh algoritma efektif. Mesin~$E$ melakukan hal ini, lalu
  berhenti dengan keluaran~$1$ ketika menemukan !!a{derivation} yang
  menunjukkan bahwa $\Proves !B$. Menurut teorema kesahihan, jika $E$
  berhenti dengan keluaran~$1$, penyebabnya ialah~$\Entails !B$. Menurut
  teorema kelengkapan, jika $\Entails !B$, terdapat !!a{derivation} yang
  menunjukkan bahwa~$\Proves !B$. Karena $E$ menghasilkan secara sistematis
  semua !!{derivation} yang mungkin, pada akhirnya mesin itu akan menemukan
  satu yang menunjukkan~$\Proves !B$, sehingga pada akhirnya berhenti dengan
  keluaran~$1$.
\end{proof}

\end{document}
